\section{Minimax precision and attainable estimation}

The precision criterion in \cref{obj:def:risk} concerns estimation of the population total treatment effect from the retained graph, treatment assignments, and realized outcomes. An attainable bound therefore begins with a score computed from those observations. Centering the binary assignments provides the treatment contrasts used in the score.

\begin{definitionv}{synth_2}[Centered treatment signs]
For a binary treatment-assignment vector \(Z\in\{0,1\}^V\), define its centered treatment-sign vector \(X\in\{-1,1\}^V\) by
\[
X_k=2Z_k-1,\qquad k\in V.
\]
Thus \(X_i=2Z_i-1\) is the own-treatment sign for recipient \(i\), and \(X_j=2Z_j-1\) is the treatment sign for source \(j\).
\end{definitionv}

For positive retention, the observable inverse-inclusion score \leanref{sym:Tq}{\ensuremath{T_q}} in \cref{obj:def:audit-score} combines each outcome with its own-treatment sign and the treatment signs of recorded sources. Weighting each recorded arrow by the reciprocal of its retention probability compensates for edge thinning under the independent audit design.

\begin{definitionv}{def:audit-score}[Inverse-inclusion score $T_q$]
\[
T_q
=
\frac2n\sum_{i\in V}Y_i(Z)
\left(X_i+q^{-1}\sum_{j:(j,i)\in H}X_j\right),
\qquad q>0.
\]
Here \(V\) is the population, \(n\) is its size, \(Z\) is the treatment-assignment vector, \(Y_i(Z)\) is unit \(i\)'s realized outcome, and \(H\) is the recorded retained-edge graph. The quantities \(X_i\) and \(X_j\) are the centered own-treatment and source-treatment signs, respectively, and \(q\) is the known edge-retention probability.
\end{definitionv}

The supplied-graph score \leanref{sym:TG}{\ensuremath{T_G}} in \cref{obj:def:oracle-score} provides the corresponding comparison when the true interference graph is available. It replaces the weighted retained-source sum with the sum over all true sources.

\begin{definitionv}{def:oracle-score}[Supplied-graph score $T_G$]
\[
T_G
=
\frac2n\sum_{i\in V}Y_i(Z)
\left(X_i+\sum_{j:(j,i)\in G}X_j\right).
\]
Here \(V\) is the population, \(n\) is its size, \(G\) is the supplied true interference graph, \(Z\) is the treatment-assignment vector, and \(Y_i(Z)\) is unit \(i\)'s realized outcome. The quantities \(X_i\) and \(X_j\) are the centered own-treatment and source-treatment signs, respectively.
\end{definitionv}

The audit decomposition in \cref{obj:lem:score-audit} separates assignment variation in the supplied-graph score from the additional variation generated by inverse-inclusion weighting. Conditional on the assignment, averaging the observable score over the audit recovers the supplied-graph score. Independent audit marks then isolate the extra variance attributable to edge collection. The resulting upper-risk envelope \leanref{sym:upper}{\ensuremath{u(n,d,q)}} in \cref{obj:def:upper-envelope} records these two components.

\begin{definitionv}{def:upper-envelope}[Upper-risk envelope $u(n,d,q)$]
\[
u(n,d,q)
=
\begin{cases}
5(d+1)^2/n+4d(1-q)/(nq),&q>0,\\
+\infty,&q=0.
\end{cases}
\]
Here \(n\) is the population size, \(d\) is the common off-diagonal in-degree and out-degree bound, and \(q\) is the known edge-retention probability.
\end{definitionv}

The first term reflects assignment variation under the common in-degree and out-degree bound. The second measures the variance cost of retaining each true arrow with probability \(q\); it decreases to zero as retention approaches one. These components give a prospective error bound in terms of population size, degree, and the collection probability.

To obtain a rule for every admissible retention probability, use the envelope to choose between a clipped score and zero. The coefficient-mass restriction in \cref{obj:ass:coefficient-mass} bounds the absolute population contrast by one, so clipping to that interval preserves the squared-error guarantee. The observable upper rule \leanref{sym:Tup}{\ensuremath{T^{\mathrm{up}}}} in \cref{obj:def:upper-rule} makes this choice explicit.

\begin{definitionv}{def:upper-rule}[Observable upper rule $T^{\mathrm{up}}$]
\[
T^{\mathrm{up}}
=
\begin{cases}
\max\{-1,\min\{1,T_q\}\},&q>0\text{ and }u(n,d,q)<1,\\
0,&\text{otherwise}.
\end{cases}
\]
Here \(T_q\) is the observable inverse-inclusion score, \(u(n,d,q)\) is the upper-risk envelope, \(n\) is the population size, \(d\) is the common off-diagonal in-degree and out-degree bound, and \(q\) is the known edge-retention probability.
\end{definitionv}

The score branch applies when its envelope improves upon the unit bound supplied by the zero rule. This construction yields the following uniform guarantee over the schedule class.

\begin{definitionv}{def:frontier-scale}[Risk scale $F(n,d,q)$]
Define the risk scale by
\[
F(n,d,q)=
\begin{cases}
\min\left\{1,\dfrac{d^2}{n[1-(1-q)^d]}\right\},&q>0,\\
1,&q=0.
\end{cases}
\]
Its domain is specified by
\begin{itemize}
\item \textbf{(Population size.)} \(n\ge4\).
\item \textbf{(Degree bound.)} \(1\le d\le n-1\).
\item \textbf{(Retention probability.)} \(q\in[0,1]\).
\end{itemize}
On this domain, \(F(n,d,q)\) is positive and finite.
\end{definitionv}

The guarantee combines variance control with the bounded range of the causal target. When assignment and audit variation are sufficiently small, the rule uses the inverse-inclusion score; when the envelope reaches one, the zero branch supplies a bounded worst-case squared loss. Thus the same observable construction covers both improving precision and saturation.

Matching lower bounds identify how this guarantee compares with the best measurable use of the complete record. The explicit risk scale \leanref{sym:F}{\ensuremath{F(n,d,q)}} in \cref{obj:def:frontier-scale} and the named attaining rule \leanref{sym:Tstar}{\ensuremath{T^\star_{n,d,q}}} in \cref{obj:def:frontier-rule} state the scale and estimator used in that characterization.

\begin{definitionv}{def:frontier-rule}[Attaining rule $T^\star_{n,d,q}$]
For every admissible \(n,d,q\), define the attaining rule on the observed record \(O\) by
\[
T^\star_{n,d,q}(O)=T^{\mathrm{up}}(O)=
\begin{cases}
\max\{-1,\min\{1,T_q(O)\}\},&q>0,\ u(n,d,q)<1,\\
0,&\text{otherwise},
\end{cases}
\]
where \(T_q\) is the observable inverse-inclusion score and \(u(n,d,q)\) is the upper-risk envelope.
\end{definitionv}

\begin{theoremv}{thm:observable-upper}[Observable upper risk bound]
Let \(V\) be a finite population with size \(n=|V|\). Suppose that:
\begin{itemize}
\item \textbf{(Population size.)} \(n\ge4\).
\item \textbf{(Degree bound.)} The integer degree bound \(d\) satisfies \(1\le d\le n-1\).
\item \textbf{(Retention probability.)} The known edge-retention probability \(q\) belongs to \([0,1]\).
\end{itemize}
Under the independent assignment--audit product design, the minimax risk \(R_{n,d}(q)\) and expected squared loss \(\mathcal L\) of \cref{obj:def:risk}, the observable rule \(T^{\mathrm{up}}\) of \cref{obj:def:upper-rule}, and the upper-risk envelope \(u(n,d,q)\) of \cref{obj:def:upper-envelope} satisfy
\[
R_{n,d}(q)
\le
\sup_{\theta\in\mathcal M_{n,d}}
\mathcal L(T^{\mathrm{up}};\theta,q)
\le
\min\{1,u(n,d,q)\},
\]
where \(\mathcal M_{n,d}\) is the fixed schedule class underlying the risk in \cref{obj:def:risk}. The rule \(T^{\mathrm{up}}\) is real-valued and measurable on the complete observed-record space.
\end{theoremv}

\begin{theoremv}{thm:observed-record-precision-frontier}[Observed-record precision frontier]
Let \(F\) be the risk-scale function in \cref{obj:def:frontier-scale}, and let the attaining rule be
\[
T^\star_{n,d,q}=T^{\mathrm{up}},
\]
where \(T^{\mathrm{up}}\) is the total observable upper rule. There exist universal real constants \(c,C_0>0\) such that, simultaneously for all parameters satisfying
\begin{itemize}
\item \textbf{(Population.)} The population size \(n\) is an integer with \(n\ge4\).
\item \textbf{(Degree.)} The degree bound \(d\) is an integer with \(1\le d\le n-1\).
\item \textbf{(Retention.)} The known retention probability \(q\) belongs to \([0,1]\).
\end{itemize}
the canonical independent assignment--audit design satisfies
\[
cF(n,d,q)
\le R_{n,d}(q)
\le
\sup_{\theta\in\mathcal M_{n,d}}
\mathcal L(T^\star_{n,d,q};\theta,q)
\le C_0F(n,d,q).
\]
Here \(\mathcal M_{n,d}\) is the fixed schedule class, \(\mathcal L(T;\theta,q)\) is expected squared estimation loss, and \(R_{n,d}(q)\) is its infimum over all measurable real-valued randomized estimators \(T(O,\xi)\), where \(O\) is the complete observed record and \(\xi\) is the independent randomization seed.

For every pair of degree and retention sequences \(d_n,q_n\) satisfying
\[
1\le d_n\le n-1,\qquad q_n\in[0,1]
\qquad\text{for every integer }n\ge4,
\]
there exists a sequence of measurable real-valued randomized estimators \(T_n(O,\xi)\) such that
\[
\sup_{\theta\in\mathcal M_{n,d_n}}
\mathcal L(T_n;\theta,q_n)\longrightarrow0
\qquad\text{as }n\longrightarrow\infty
\]
if and only if
\[
\frac{d_n^2}{n}\longrightarrow0
\qquad\text{and}\qquad
\begin{cases}
+\infty,&q_n=0,\\[2pt]
\dfrac{d_n}{nq_n},&q_n>0
\end{cases}
\longrightarrow0.
\]
The second convergence is in the extended nonnegative real numbers.

For every integer \(n\ge4\) and \(1\le d\le n-1\), the endpoint scales are
\[
F(n,d,0)=1,
\qquad
F(n,d,1)=\min\left\{1,\frac{d^2}{n}\right\}.
\]
Moreover, there exist universal real constants \(c_1,C_1>0\) such that, simultaneously for every integer \(n\ge4\), every integer \(1\le d\le n-1\), and every \(0<q\le1\),
\[
c_1\min\left\{1,\frac{d^2}{n}+\frac{d}{nq}\right\}
\le F(n,d,q)
\le
C_1\min\left\{1,\frac{d^2}{n}+\frac{d}{nq}\right\}.
\]
\end{theoremv}

The characterization separates two sources of imprecision: the degree component \(d^2/n\) and the edge-collection component \(d/(nq)\). The universal lower and upper constants \leanref{sym:c}{\ensuremath{c}} and \leanref{sym:C0}{\ensuremath{C_0}} apply simultaneously across the stated parameter domain. Consequently, the observable rule attains the minimax scale even when competing estimators use all assignment and outcome coordinates and independent randomization.

The factor \leanref{sym:p}{\ensuremath{1-(1-q)^d}} is the probability that at least one of \(d\) independently audited arrows is retained. Its role in the scale describes the collection transition: when \(dq\) is small, edge collection determines the leading precision cost; once \(dq\) is of constant order or larger, the degree component determines the scale up to universal factors. The transition occurs around retention of order \(1/d\). The minimum with one expresses saturation of worst-case squared risk: whenever the uncapped scale is at least one, minimax risk is bounded above and below by positive universal constants.

For admissible degree sequences \leanref{sym:dseq}{\ensuremath{d_n}} and retention sequences \leanref{sym:qseq}{\ensuremath{q_n}}, the consistency statement gives an exact joint requirement. Degree must grow more slowly than the square root of population size, and retention must be positive eventually with \(nq_n/d_n\) diverging. Under these conditions, the named attaining rules have uniformly vanishing expected squared loss. Necessity follows from the lower bound over all measurable randomized estimators, so these conditions characterize the estimation problem under the declared design and schedule class.

The converse uses blocks of recipients that share response values while their source allocations are only partly revealed by the audit. Before treatment and audit draws, each block receives a fixed but unknown baseline and a random allocation of sources. The observed outcomes can therefore teach the estimator about hidden associations, so the comparison retains repeated responses and conditions the likelihood on the observed global treated-source count. Translation control, Walsh energy, and hidden-allocation contraction show that opposite-sign block schedules remain difficult to distinguish even with that outcome-driven learning; see \cref{obj:lem:block-law,obj:lem:baseline-translation-affinity,obj:lem:walsh-channel-energy,obj:lem:hidden-allocation-contraction,obj:thm:block-testing-scale}. The two-prior reduction in \cref{obj:lem:full-record-two-prior-risk} then turns this complete-record testing bound into the minimax lower bound. Endpoint and explicit-constant deductions are collected in \cref{obj:thm:supported-boundaries,obj:thm:frontier-answer}.
